\section{The primitive-core boundary and the next gate}
\label{sec:tpc29-primitive-boundary}

The two theorems remove large shared content on both sides of the
fiber-length transition.  They do not remove the whole complement.
This section records the exact remainder and prevents a sign-blind
interpretation of the result.

\subsection{The sparse method boundary}

In exponent notation, a primitive complement pair has $c=0$.  The
sparse condition \eqref{eq:tpc29-sparse-exponent} becomes
\begin{equation}
 a+b\ge\jmath,
 \label{eq:tpc29-primitive-sparse-condition}
\end{equation}
whereas the incidence saving is
\begin{equation}
 \eta_{\rm sp}(a,b,0)
 =(\jmath-a-b)_+=0.
 \label{eq:tpc29-primitive-zero-saving}
\end{equation}
Thus \cref{thm:tpc29-weighted-block} reaches only the natural scale on
the coprime sparse core.  This is a sharp boundary of the present
absolute-incidence argument, not a lower bound for the signed
arithmetic sum.

For a fixed $\eta>0$, after removing
\cref{thm:tpc29-rich-sparse}, the remaining sparse region may be
written as
\begin{equation}
 [r,s]\ge J,
 \qquad
 \frac{[r,s]}{(r,s)}>JX^{-\eta}.
 \label{eq:tpc29-remaining-sparse}
\end{equation}
It contains every coprime pair with $rs\ge J$.  Although each such
fiber has $O(1)$ orbit points, the family has enough modulus pairs to
restore the natural scale.  A short fiber by itself is therefore not a
cancellation theorem.

\subsection{The remaining long fibers}

For $[r,s]\le J$, \cref{thm:tpc29-long-content} removes
$(r,s)>X^\eta$ for any fixed $\eta>0$.  What remains has
\begin{equation}
 [r,s]\le J,\qquad (r,s)\le X^\eta.
 \label{eq:tpc29-remaining-long}
\end{equation}
On any further subpacket carrying the additional generic row
separation
$|m-n|\ge QX^{-\eta}$, the affine determinant in
\eqref{eq:tpc29-affine-determinant} satisfies
\begin{equation}
 \left|\frac{h(m-n)}{(r,s)}\right|
 \gg_{h}QX^{-2\eta}.
 \label{eq:tpc29-large-remaining-determinant}
\end{equation}
Thus the separated subpacket is not a small-determinant degeneracy.
This conditional observation is not used in either closure theorem
and does not remove any unseparated row pairs.  The remaining long
problem consists of small-content, polynomial-slope one- and
two-point reflected M\"obius families: one quotient M\"obius factor
moves in a mixed channel, and two move in the both-ultra channel.

For a fixed compatible quadruple $(m,n,r,s)$ in the both-ultra
channel, the logarithmically averaged two-point Elliott theorem is a
fixed-form cancellation benchmark for the two affine quotient forms
in \eqref{eq:tpc29-U-form}--\eqref{eq:tpc29-V-form}
\citep{Tao2016}.  A mixed channel has only a one-point moving M\"obius
factor.  Neither fixed-form benchmark supplies an effective rate
uniform over the polynomially growing slopes, the divisor family, and
the physical row weights.  It therefore cannot simply be summed over
\eqref{eq:tpc29-remaining-long}.

\subsection{A sufficient next interface}

The next useful statement is a family estimate, not a pointwise
asymptotic for an $O(1)$-point fiber.  One possible interface is the
following.

\begin{definition}[Primitive reflected-family dispersion]
\label{def:tpc29-next-interface}
For a fixed content threshold $X^\eta$, let $\cP_\eta$ be the union of
the remaining regions \eqref{eq:tpc29-remaining-sparse} and
\eqref{eq:tpc29-remaining-long}, with the actual mixed and both-ultra
coefficient types retained.  A primitive reflected-family dispersion
estimate is a bound
\begin{equation}
 \sum_{\textup{three raw channels}}
 \cS(\cP_\eta)
 \ll_{B,\eta,h,\mathscr D}XQ(\log X)^{-B}
 \label{eq:tpc29-next-dispersion}
\end{equation}
for every fixed $B>0$, uniformly over the physical opened-row weights
and row-only mask.
\end{definition}

Equation \eqref{eq:tpc29-next-dispersion} is not assumed or proved in
this paper.  Its formulation preserves the joint cutoff difference:
one must not first assert that each complete mixed channel or the
complete both-ultra channel is negligible.  The content-rich pieces
may be removed separately only because the present theorems bound
them absolutely.

\subsection{Status after this paper}

\begin{center}
\begin{tabular}{>{\raggedright\arraybackslash}p{0.43\textwidth}
                >{\raggedright\arraybackslash}p{0.46\textwidth}}
\toprule
Object & Status \\
\midrule
Selected calibrated square through $T$
& Closed in TPC-28 under its physical conductor interfaces. \\
\addlinespace
Two row-drift legs beyond $T$
& Fixed-power small by \cref{prop:tpc29-drift-small}. \\
\addlinespace
Long fibers $[r,s]\le J$ with $(r,s)>X^\eta$
& Closed with fixed-power saving by
  \cref{thm:tpc29-long-content}. \\
\addlinespace
Sparse fibers with $[r,s]\ge J$ and
$[r,s]/(r,s)\le JX^{-\eta}$
& Closed with fixed-power saving by
  \cref{thm:tpc29-rich-sparse}. \\
\addlinespace
Small-content long fibers and primitive sparse fibers
& Open; require signed family dispersion such as
  \eqref{eq:tpc29-next-dispersion}. \\
\addlinespace
Complete TPC-18 residual and endpoint reassembly
& Not proved. \\
\addlinespace
Positivity, Hardy--Littlewood asymptotic, twin primes, or general
parity breakthrough
& Not claimed and not implied. \\
\bottomrule
\end{tabular}
\end{center}

This is a genuine reduction of the open complement geometry: the
remaining long core has small gcd, whereas the remaining sparse core
has large primitive lcm $[r,s]/(r,s)$.  Neither remaining condition is
a signed cancellation theorem, and this is not a solution of the
parity-sensitive core.
